% !TeX root = ../main.tex

\section{Progettazione Fisica}
La progettazione fisica rappresenta la traduzione concreta dello schema logico nel DBMS scelto, nel nostro caso \textbf{PostgreSQL}. In questa fase si definiscono le tabelle effettive, i tipi di dato, i vincoli di integrità esprimibili direttamente nel DDL e gli indici necessari a supportare in modo efficiente le operazioni più frequenti previste dal progetto.

Dal punto di vista operativo, l'infrastruttura del database è stata organizzata mediante una fase iniziale di creazione e reinizializzazione del database, seguita dalla definizione dello schema applicativo \texttt{consorzio}. Tutte le tabelle, i vincoli e la logica successiva vengono infatti collocati all'interno di tale schema, così da mantenere ordinata e separata l'intera base di dati.

\subsection{Scelte implementative e organizzazione dello schema}
La realizzazione fisica è stata organizzata attorno allo schema \texttt{consorzio}, impostato come \texttt{search\_path} di riferimento per gli script SQL. L'inizializzazione prevede la chiusura delle connessioni attive al database, la sua ricreazione e la definizione dello schema applicativo.

\begin{lstlisting}[caption={Inizializzazione del database e dello schema}]
SELECT pg_terminate_backend(pg_stat_activity.pid)
FROM pg_stat_activity
WHERE pg_stat_activity.datname = '${DB_NAME}' AND pid <> pg_backend_pid();

DROP DATABASE IF EXISTS ${DB_NAME};
CREATE DATABASE ${DB_NAME};

create schema if not exists consorzio;
set search_path to consorzio;
\end{lstlisting}

Per quanto riguarda la modellazione dei domini, si è scelto di sfruttare i tipi enumerativi di PostgreSQL per gli attributi che ammettono un insieme chiuso di valori. In particolare, il tipo di interfaccia e il tipo di operazione sono stati modellati tramite \texttt{ENUM}, evitando così l'uso ripetuto di vincoli \texttt{CHECK} sulle stesse colonne e rendendo il dominio dei valori più esplicito a livello di schema.

\begin{lstlisting}[caption={Definizione dei tipi enumerativi}]
create type tipo_interfaccia as enum ('cassa', 'sportello');
create type tipo_operazione as enum ('prelievo', 'deposito', 'saldo', 'lista_movimenti');
\end{lstlisting}

Dal punto di vista dei tipi di dato:
\begin{itemize}
    \item gli identificativi principali del sistema bancario sono stati realizzati prevalentemente come interi autogenerati tramite \texttt{GENERATED ALWAYS AS IDENTITY};
    \item gli attributi descrittivi, come nomi e indirizzi, sono stati modellati con \texttt{VARCHAR};
    \item gli importi monetari, i limiti di spesa e le disponibilità sono stati rappresentati tramite \texttt{NUMERIC(x, 2)}, così da evitare approssimazioni indesiderate mantenendo due cifre dopo la virgola;
    \item le informazioni temporali sono state memorizzate con i tipi \texttt{DATE} e \texttt{TIME}, coerentemente con il modello logico.
\end{itemize}

Queste scelte permettono di mantenere lo schema aderente al modello progettato e, allo stesso tempo, di sfruttare in modo appropriato le caratteristiche offerte da PostgreSQL.

\subsection{Definizione delle tabelle}
Le tabelle sono state implementate seguendo la struttura dello schema logico, distinguendo tra organizzazione bancaria, soggetti del sistema, conti e strumenti di pagamento, interfacce operative e registrazione delle operazioni.

\subsubsection{Organizzazione bancaria}
Le prime due tabelle descrivono la struttura del consorzio bancario: \texttt{banca} rappresenta gli istituti aderenti, mentre \texttt{filiale} modella le sedi operative appartenenti alle singole banche.

\begin{lstlisting}[caption={Tabelle banca e filiale}]
create table banca (
    codiceBanca integer generated always as identity primary key,
    nomeBanca varchar(50) not null,
    commissioneAltri numeric(4, 2) not null,
    commissioneConsorzio numeric(4, 2) not null
);

create table filiale (
    codiceFiliale integer generated always as identity primary key,
    indirizzo varchar(100) not null,
    bancaReferente integer not null,
    foreign key (bancaReferente) references banca(codiceBanca)
);
\end{lstlisting}

In questa fase le commissioni sono memorizzate direttamente nella tabella \texttt{banca}, poiché costituiscono un dato strutturale necessario al successivo calcolo dei costi applicati ai prelievi.

\subsubsection{Soggetti del sistema}
Per rappresentare correttamente il fatto che una stessa persona possa essere cliente, impiegato, oppure entrambe le cose, la generalizzazione \texttt{Persona} -- \texttt{Cliente}/\texttt{Impiegato} è stata mantenuta mediante tabelle distinte collegate tra loro tramite chiavi esterne.

\begin{lstlisting}[caption={Tabelle persona, cliente e impiegato}]
create table persona (
    codiceFiscale varchar(16) primary key,
    indirizzo varchar(100) not null,
    nome varchar(30) not null
);

create table cliente (
    codiceFiscale_c varchar(16) primary key
        references persona(codiceFiscale) ON DELETE CASCADE
);

create table impiegato (
    codiceFiscale_i varchar(16) primary key
        references persona(codiceFiscale) ON DELETE CASCADE,
    idImpiegato integer generated always as identity unique,
    dataAssunzione date not null
);
\end{lstlisting}

L'attributo \texttt{idImpiegato} è stato introdotto come identificatore operativo univoco dell'impiegato. Esso non sostituisce il codice fiscale come chiave primaria della tabella, ma viene utilizzato nelle relazioni operative e nei trigger, risultando più pratico nella gestione interna dei trasferimenti e delle assegnazioni.

La storicizzazione del servizio prestato in filiale è stata implementata tramite la tabella \texttt{presta\_servizio\_in}.

\begin{lstlisting}[caption={Tabella presta\_servizio\_in}]
create table presta_servizio_in (
    idImpiegato_psi integer,
    codiceFiliale_psi integer,
    dataInizio date not null,
    dataFine date,
    primary key (idImpiegato_psi, codiceFiliale_psi),
    foreign key (idImpiegato_psi) references impiegato(idImpiegato) ON DELETE CASCADE,
    foreign key (codiceFiliale_psi) references filiale(codiceFiliale) ON DELETE CASCADE
);
\end{lstlisting}

La tabella registra il legame tra impiegato e filiale, insieme alle date di inizio e fine del servizio. La chiave primaria composta impedisce che la stessa coppia impiegato-filiale venga registrata più di una volta, mentre i vincoli temporali più specifici vengono demandati a controlli aggiuntivi.

\subsubsection{Conti e carte Bancomat}
Il nucleo bancario del progetto è rappresentato dalle tabelle \texttt{conto}, \texttt{intestato\_a} e \texttt{cartaBancomat}. La prima descrive il conto corrente, la seconda realizza la cointestazione tra clienti e conti, mentre la terza modella le carte Bancomat associate a un cliente e a un conto.

\begin{lstlisting}[caption={Tabelle conto e intestato\_a}]
create table conto (
    numeroConto integer generated always as identity primary key,
    dataApertura date not null,
    dataEstinzione date,
    ammontare numeric(17,2),
    filialeReferente integer not null,
    foreign key (filialeReferente) references filiale(codiceFiliale)
);

create table intestato_a (
    codiceFiscale_int varchar(16),
    numeroConto_int integer,
    primary key (codiceFiscale_int, numeroConto_int),
    foreign key (codiceFiscale_int) references cliente(codiceFiscale_c) ON DELETE CASCADE,
    foreign key (numeroConto_int) references conto(numeroConto) ON DELETE CASCADE
);
\end{lstlisting}

\begin{lstlisting}[caption={Tabella cartaBancomat}]
create table cartaBancomat (
    passwordBancomat varchar(25) primary key,
    numeroConto_cb integer not null,
    codiceFiscale_cb varchar(16) not null,

    limiteSpesaGiornaliero numeric(10, 2) not null,
    limiteSpesaMensile numeric(14, 2) not null,

    spesaGiornaliera numeric(10, 2) default 0 not null,
    spesaMensile numeric(14, 2) default 0 not null,

    foreign key (numeroConto_cb) references conto(numeroConto) ON DELETE CASCADE,
    foreign key (codiceFiscale_cb) references cliente(codiceFiscale_c) ON DELETE CASCADE
);
\end{lstlisting}

Nella tabella \texttt{cartaBancomat} sono presenti anche gli attributi \texttt{spesaGiornaliera} e \texttt{spesaMensile}, mantenuti direttamente nel database per consentire controlli rapidi sui limiti operativi senza dover ricalcolare ogni volta lo storico delle operazioni effettuate.

\subsubsection{Interfacce e operazioni}
Le casse e gli sportelli automatici, inizialmente distinti a livello concettuale, sono stati ricondotti all'unica tabella \texttt{interfaccia}, differenziata tramite l'attributo \texttt{tipo}. Allo stesso modo, le operazioni bancarie sono state accorpate nella tabella \texttt{operazione}, con un attributo discriminante anch'esso tipizzato.

\begin{lstlisting}[caption={Tabelle interfaccia e operazione}]
create table interfaccia (
    idInterfaccia integer generated always as identity primary key,
    tipo tipo_interfaccia not null,
    disponibilita numeric(10, 2) not null,
    numeroOperazioni integer not null default 0,
    codiceFiliale_i integer not null,
    foreign key (codiceFiliale_i) references filiale(codiceFiliale)
);

create table operazione (
    idOperazione integer generated always as identity primary key,
    dataOperazione date default CURRENT_DATE not null,
    oraOperazione time default CURRENT_TIME not null,
    tipo tipo_operazione not null,
    ammontare numeric(10,2),
    numeroConto_o integer not null,
    idInterfaccia_o integer not null,
    passwordBancomat_o varchar(25) not null,
    foreign key (numeroConto_o) references conto(numeroConto),
    foreign key (idInterfaccia_o) references interfaccia(idInterfaccia),
    foreign key (passwordBancomat_o) references cartaBancomat(passwordBancomat)

  ON DELETE CASCADE
);
\end{lstlisting}

L'uso dei valori di default per \texttt{dataOperazione} e \texttt{oraOperazione} permette di registrare automaticamente il momento dell'operazione se non specificato esplicitamente. Inoltre, l'attributo \texttt{ammontare} è lasciato opzionale, poiché operazioni come saldo e lista movimenti non richiedono un importo numerico associato.

\subsection{Vincoli di integrità esprimibili nello schema fisico}
Una parte dei vincoli progettuali è stata implementata direttamente nel DDL mediante chiavi, riferimenti esterni, tipi enumerativi e vincoli aggiuntivi introdotti successivamente con \texttt{ALTER TABLE}.

Oltre ai vincoli già presenti nelle definizioni delle tabelle, nel file \texttt{02\_constraints.sql} sono stati aggiunti due controlli espliciti:
\begin{itemize}
    \item un vincolo di non negatività sull'ammontare delle operazioni;
    \item un vincolo di coerenza temporale sulle date di \texttt{presta\_servizio\_in}.
\end{itemize}

\begin{lstlisting}[caption={Vincoli aggiuntivi sullo schema fisico}]
alter table operazione
    add constraint chk_ammontare_positivo
    check (ammontare >= 0);

alter table presta_servizio_in
    add constraint chk_date
    check (dataFine is null or dataFine >= dataInizio);
\end{lstlisting}

Questi vincoli impediscono già a livello strutturale:
\begin{itemize}
    \item la registrazione di operazioni con importo negativo;
    \item l'inserimento di intervalli temporali incoerenti nel servizio degli impiegati.
\end{itemize}

È importante osservare che gli \texttt{ENUM} svolgono anch'essi una funzione di integrità, poiché impediscono l'inserimento di valori arbitrari nei campi \texttt{tipo} di \texttt{interfaccia} e \texttt{operazione}.

\subsection{Indici e ottimizzazione delle interrogazioni}
Accanto alla definizione delle tabelle e dei vincoli, la progettazione fisica ha incluso anche la creazione di indici su attributi particolarmente rilevanti per le interrogazioni previste dal progetto. La scelta è stata effettuata in coerenza con il carico applicativo descritto nelle sezioni precedenti e con le query effettivamente implementate.

Sono stati introdotti indici sulle principali chiavi esterne, così da velocizzare join e attraversamenti frequenti tra le tabelle.

\begin{lstlisting}[caption={Indici sulle principali chiavi esterne}]
create index idx_filiale_banca_fk on filiale(bancaReferente);
create index idx_conto_filiale_fk on conto(filialeReferente);
create index idx_interfaccia_filiale_fk on interfaccia(codiceFiliale_i);
create index idx_operazione_conto_fk on operazione(numeroConto_o);
create index idx_operazione_interfaccia on operazione(idInterfaccia_o);
create index idx_intestato_conto_fk on intestato_a(numeroConto_int);
\end{lstlisting}

Oltre a questi, sono stati definiti anche indici specifici per alcune query più selettive:

\begin{lstlisting}[caption={Indici specializzati e parziali}]
create index idx_conto_attivo_data
    on conto(dataApertura)
    where dataEstinzione is null;

create index idx_operazione_ammontare_val
    on operazione(ammontare);

create index idx_servizio_attivo
    on presta_servizio_in(idImpiegato_psi)
    where dataFine is null;
\end{lstlisting}

In particolare:
\begin{itemize}
    \item l'indice parziale su \texttt{conto(dataApertura)} considera soltanto i conti attivi, ottimizzando le interrogazioni che lavorano sui rapporti non estinti;
    \item l'indice su \texttt{operazione(ammontare)} supporta le query che filtrano o aggregano in base all'importo dell'operazione;
    \item l'indice parziale su \texttt{presta\_servizio\_in(idImpiegato\_psi)} accelera le interrogazioni sugli impiegati attualmente in servizio, limitando l'indicizzazione alle sole tuple con \texttt{dataFine} nulla.
\end{itemize}

Queste scelte permettono di migliorare le prestazioni senza introdurre ridondanze aggiuntive nel modello dati.

\subsection{Vincoli non esprimibili mediante solo DDL}
Non tutti i vincoli richiesti dalla traccia possono essere espressi unicamente mediante \texttt{PRIMARY KEY}, \texttt{FOREIGN KEY}, \texttt{CHECK} o tipi enumerativi. Alcune regole di business dipendono infatti da confronti tra più tabelle, da verifiche sullo stato corrente del database oppure da aggiornamenti automatici conseguenti a una operazione.

Negli script effettivamente adottati, tali aspetti sono demandati a funzioni e trigger \texttt{PL/pgSQL}, che verranno descritti nella fase di implementazione. In particolare, la logica procedurale si occupa di:
\begin{itemize}
    \item verificare che la carta Bancomat usata in una operazione sia effettivamente associata al conto su cui si sta operando;
    \item impedire trasferimenti non validi degli impiegati, sia in caso di ritorno nella stessa filiale sia in caso di permanenza inferiore a un mese;
    \item calcolare le commissioni sui prelievi in base alla banca del conto e alla banca dell'interfaccia utilizzata;
    \item bloccare i prelievi che eccedono il saldo del conto, la disponibilità dell'interfaccia o i limiti giornalieri e mensili della carta;
    \item aggiornare automaticamente, dopo ogni operazione, saldo del conto, disponibilità dell'interfaccia, contatore delle operazioni e spese accumulate sulla carta;
    \item impedire operazioni su conti già estinti.
\end{itemize}

Di conseguenza, la progettazione fisica non si limita alla sola definizione delle tabelle, ma prepara anche il terreno alla logica attiva del database, che completa i vincoli statici con controlli dinamici e comportamenti automatici.

Nel complesso, lo schema fisico realizzato in PostgreSQL rimane fedele alla progettazione logica, ma introduce gli strumenti concreti necessari a rendere il sistema eseguibile, coerente e sufficientemente efficiente anche su basi di dati di dimensione non trascurabile.
